% OWD 2023
% Placement Interview Presentation for 'owd-euses'
%
% To generate the man page PDF in the listings/ directory:
%   $ man -t manual.1 > manual.ps
%   $ ps2pdf manual.ps

\documentclass{beamer}

% BEGIN COMPANY INFORMATION

\newcommand\companyname{Thales}
\date{27th September 2023} % Date of the placement interview

% END COMPANY INFORMATION

\usepackage{minted, calc, caption}

\usetheme{CambridgeUS}
\setbeamertemplate{caption}[numbered]
\numberwithin{figure}{section}
\beamerdefaultoverlayspecification{<+->}
\titlegraphic{\includegraphics[width=.25\textwidth]{./gentoo-horizontal}}

\newcommand\chext{.{\fontfamily{cmr}\{}c,h{\fontfamily{cmr}\}}}
\newcommand\printmanpage[1]{\fbox{%
    \includegraphics[page=#1, width=.3\linewidth]{listings/manual.pdf}}}

% Note to self: is there a nicer way to do this without an excessive amount of
% work? There doesn't seem to be any immediately obvious way of knowing how many
% slides are required for any given source code file. Perhaps minted exposes
% some tests in a publicly accessible interface?
\newcounter{firstsrcline}
\newcounter{lastsrcline}
\newcommand\srclinelimit{29}
\newcommand\resetsrclines{%
    \setcounter{firstsrcline}{1}
    \setcounter{lastsrcline}{\thefirstsrcline+\srclinelimit}
}
\newcommand\continuesource[1]{%
    \inputminted[firstline=\thefirstsrcline, lastline=\thelastsrcline,
        firstnumber=\thefirstsrcline]{c}{listings/#1}
    \setcounter{firstsrcline}{\thelastsrcline+1}
    \setcounter{lastsrcline}{\thefirstsrcline+\srclinelimit}
}
\renewcommand\theFancyVerbLine{\color{gray}\slshape\arabic{FancyVerbLine}}
\setminted{
    fontsize=\tiny,
    style=vs,
    linenos,
    numbersep=8pt,
    xleftmargin=12pt,
}

\hypersetup{
    colorlinks,
    allcolors=.,
    urlcolor=gray,
}

\newcommand\programname{\raisebox{.5pt}{\ttfamily owd-euses}}
\title[Placement Interview Presentation]{The \programname\ Querying Tool}
\author[Oliver Dixon]{
    \texorpdfstring{Oliver Dixon \raisebox{.5pt}{%
        \texttt{\footnotesize\href{mailto:Oliver Dixon <od641@york.ac.uk>}
            {<od641@york.ac.uk>}}
    }}{Oliver Dixon}}
\institute[University of York]{Department of Computer Science, University of
    York \and Department of Mathematics, University of York}
\subtitle{\companyname\ Placement Interview Presentation}

\begin{document}
\frame{\titlepage}
\section{Introduction, Motivation, and Project Specification}
\begin{frame}
    \frametitle{Overview of the Environment: Gentoo Linux}
    \begin{itemize}
        \item Gentoo Linux is a source-based distribution of the Linux operating
            system; such distributions typically include package managers that
            encourage, or mandate, the local compilation of most packages. This
            stands in contrast to the previously established model of
            downloading and installing pre-built binaries.
        \item A source-based distribution model has a number of advantages for
            power users, namely the ability to enable and disable compile-time
            features as desired, resulting in a more compact and suitably
            equipped binary.
        \item With such extensive customisability, a powerful querying system is
            required to determine the various features that are supported by
            each of the available packages.
    \end{itemize}
\end{frame}
\begin{frame}
    \frametitle{Existing Solutions}
    Due to the age and popularity of Gentoo Linux, a multitude of querying tools
    exist for this purpose:
    \pause

    \begin{itemize}
        \item \emph{eix}: very powerful and well-established, but requires
            building and maintaining a centralised binary cache of every
            installable package.
        \item \emph{euses}: one of the original cornerstones of the Gentoo
            toolset, but contains bugs, error-prone and non-portable code, and
            is incompatible with newer versions of Gentoo.
        \item \emph{portageq}: the standard Gentoo querying tool; fully
            compliant, but proves to be very non-performant on large software
            repositories due to its antiquated Python foundations.
    \end{itemize}
\end{frame}
\begin{frame}
    \frametitle{Project Requirements}
    \programname\ was devised in an effort to unify the advantages of these
    existing solutions while eliminating the problems relating to safety,
    incompatibility, and performance-preserving scalability. The final tool must
    possess certain critical attributes:
    \pause

    \begin{enumerate}
        \item be fully standards-compliant with the Gentoo \emph{Package Manager
            Specification} (PMS) and sufficiently customisable to work on future
            revisions;
        \item operate on existing (plain text) files in the established
            directory structures, instead of generating and periodically
            maintaining binary caches;
        \item seamlessly scale on very large software repositories;
        \item be \emph{correct}: undergoing rigorous unit-testing,
            stress-testing under hostile environments, and never entering an
            undefined or non-deterministic state.
    \end{enumerate}
\end{frame}
\section{Technical Overview}
\begin{frame}
    \frametitle{Details of the Implementation and Source Structure}
    Written in C99, \programname\ was intended to run on POSIX-compliant
    Gentoo/PMS-conforming UNIX-like systems. Given the four major design
    principles, the codebase is very compact, currently standing at slightly
    over 2000 lines. The source is clearly sectioned into distinct functions and
    well-abstracted with regard to near-complete separation of code and data.
    \pause

    \begin{center}
        \begin{tabular}{ll}
            \hline
            \textbf{Translation Unit(s)} & \textbf{Purpose} \\
            \hline
            \texttt{euses\chext} & Provide the core parsing and searching
                routines \\
            \texttt{globbing\chext} & Utilise POSIX regex support \\
            \texttt{args\chext} & Parse and manage command-line arguments \\
            \texttt{stack\chext} & Implement the standard stack ADT \\
            \texttt{converse\chext} & Manage status- and error-reporting \\
            \texttt{colour.h} & Provide helpers for ANSI text-colouring \\
            \hline
        \end{tabular}
    \end{center}
\end{frame}
\begin{frame}
    \frametitle{Source Preview: Stack Interface (1)}
    \resetsrclines
    \continuesource{stack.h}
\end{frame}
\begin{frame}
    \frametitle{Source Preview: Stack Interface (2)}
    \continuesource{stack.h}
\end{frame}
\begin{frame}
    \frametitle{Source Preview: Stack Implementation (1)}
    \resetsrclines
    \continuesource{stack.c}
\end{frame}
\begin{frame}
    \frametitle{Source Preview: Stack Implementation (2)}
    \continuesource{stack.c}
\end{frame}
\begin{frame}
    \frametitle{Source Preview: Stack Implementation (3)}
    \continuesource{stack.c}
\end{frame}
\begin{frame}
    \frametitle{Abridged Source Preview: Argument Interface (1)}
    \resetsrclines
    \continuesource{args.h}
\end{frame}
\begin{frame}
    \frametitle{Abridged Source Preview: Argument Interface (2)}
    \continuesource{args.h}
\end{frame}
\begin{frame}
    \frametitle{Abridged Source Preview: Argument Interface (3)}
    \continuesource{args.h}
\end{frame}
\begin{frame}
    \frametitle{Abridged Source Preview: Argument Interface (4)}
    \continuesource{args.h}
\end{frame}
\begin{frame}
    \frametitle{Source Preview: Two-Way Heuristic Searching Driver}
    \resetsrclines
    \continuesource{twoway.c}
\end{frame}
\begin{frame}
    \frametitle{Example Usage: List Repositories and Perform a Query}
    Invocation: \texttt{owd-euses --repo-paths --verbose --list-repos qt5}
    \par\vfill
    \inputminted{text}{listings/qt.example}
    \vfill\par
    Equivalent Invocation with Short Arguments: \texttt{owd-euses -pvr qt5}
\end{frame}
\begin{frame}
    \frametitle{Example Usage: Multiple Global Queries with Identification over
        a Nonstandard Directory Structure}
    Invocation: \texttt{PORTAGE\_CONFIGROOT=/etc/portage owd-euses --global
        --repo-names --no-case --print-needles tls ssl vpn}
    \par\vfill
    \inputminted{text}{listings/tls.example}
    \vfill\par
    Equivalent Invocation with Short Arguments:
    \texttt{PORTAGE\_CONFIGROOT=/etc/portage owd-euses -gnce tls ssl vpn}
\end{frame}
\begin{frame}
    \frametitle{Example Usage: Informative Error-Handling}
    Invocation: \texttt{PORTAGE\_CONFIGROOT=/var/db/repos/gentoo owd-euses}
    \par Problem: The given directory structure is non-conforming.
    \par\vfill\inputminted{text}{listings/error-dir.example}

    \vfill
    Invocation: \texttt{owd-euses -gk qt5}
    \par Problem: The search-space specification is contradictory.
    \par\vfill\inputminted{text}{listings/error-args.example}
\end{frame}
\section{Release and Reception}
\begin{frame}
    \frametitle{Packaging, Distribution, and Maintenance}
    \programname\ was exceptionally well-received by the community, and
    undoubtedly achieved its ambitions.
    \pause

    \begin{itemize}
        \item The package was shipped with makefiles and Portage-compatible
            installation scripts (\emph{ebuilds}); this allowed for easy
            adoption by the wider Gentoo community.
        \item After first being advertised on the Gentoo developers' mailing
            list, \programname\ was added to the official software repository,
            being mirrored by hundreds of servers throughout the world.
        \item The consistently rigorous Doxygen documentation combined with the
            \texttt{roff}-compatible manual page guarantees a pleasant
            experience for developers and end-users.
        \item A central toolkit for similar tasks has since been released by the
            central Gentoo Infrastructure Team, but \programname\ remains in use
            to a lesser extent.
    \end{itemize}
\end{frame}
% \section{Future Work}
%
% Note to self: keep this out of the presentation; the interviewers will likely
% ask follow-up questions regarding potential improvements. Such questions can
% be easily answered: add regex support for match queries, add wide character
% support, threading the search algorithms, optimising for shorter needle
% lengths, etc.
\section{Questions}
\begin{frame}
    \frametitle{Any questions?}
    \begin{figure}[h!]
        \centering
        \printmanpage{1}
        \printmanpage{2}
        \printmanpage{3}
        \vspace{.5em}
        \caption*{\slshape The \programname\ Manual Page}
    \end{figure}
\end{frame}
\end{document}

